% ============ 第2章 §2.16 极值问题的解(原书 pp.165-173) ============
\section{极值问题的解}\label{sec:2.16}

当 $1 \leq p \leq \infty$ 时,对 $\phi \in (H^{p})^{*}$,其极值问题的解中满足 $\phi(f) > 0$ 者的全体记为 $S_{\phi}^{p}$.即
\begin{equation*}
S_{\phi}^{p} = \{f \in S^{p}; \phi(f) = \|\phi\|\},
\end{equation*}
其中 $S^{p} = \{f \in H^{p}; \|f\|_{p} = 1\}$.$S_{\phi}^{p}$ 是凸集合.设 $\langle S_{\phi}^{p} \rangle$ 为包含 $S_{\phi}^{p}$ 的 $H^{p}$ 的最小线性子空间.由定理~\ref{thm:2.14.4},当 $1 < p < \infty$ 时 $\langle S_{\phi}^{p} \rangle$ 是 1 维的.又当 $p = \infty$ 且 $\phi \in L^{1}/zH^{1}$ 时 $\langle S_{\phi}^{p} \rangle$ 是 1 维的.但当 p = 1 时,即使 $\phi \in C/zA$,$\langle S_{\phi}^{p} \rangle$ 也未必是 1 维的.当 $1/p + 1/q = 1$ 时,对 $k \in L^{q}$,
\begin{equation*}
\phi(f) = \frac{1}{2\pi i}\int_{|z|=1} f(z) k(z)\dd z \quad (f \in H^{p})
\end{equation*}
时记 $S_{\phi}^{p} = S_{k}^{p}$.当 $p \neq \infty$ 时,由定理~\ref{thm:2.13.1},对 $\phi \in (H^{p})^{*}$ 总存在上述那样的 k.

引理~\ref{lem:2.14.3} 给出了 $\phi \in (H^{p})^{*}$ 的极值核与极值解的关系,由此立刻导出下面的命题.

\begin{proposition}\label{prop:2.16.1}
设 $1 \leq p < \infty$,$\phi \in (H^{p})^{*}$ 且 $\|\phi\| = 1$.设 K 为 $\phi$ 的极值核,则
\begin{equation*}
S_{\phi}^{p} = \left\{f \in S^{p}; K(e^{i\theta}) = e^{-i\theta}\frac{|f(e^{i\theta})|^{p}}{f(e^{i\theta})} \text{ a.e. } \theta \right\}.
\end{equation*}
这里,p = 1 时假定 $S_{\phi}^{1} \neq \emptyset$.
\end{proposition}

\begin{proof}
由引理~\ref{lem:2.14.3},若 $f \in S_{\phi}^{p}$,则 $e^{i\theta}K(e^{i\theta})f(e^{i\theta}) = |K(e^{i\theta})||f(e^{i\theta})|$ 且 $|K(e^{i\theta})| = |f(e^{i\theta})|^{p/q}$,由 $1/p + 1/q = 1$ 即明.
\end{proof}

k 的具体例子如下.

\begin{example}\label{exa:2.16.2}
设 k 为 $\phi$ 的积分表示的核.设 $f(z) = \sum_{j=0}^{\infty} a_{j}z^{j} \in H(D)$.

(1) 若 $k(z) = \sum_{j=1}^{n}\frac{c_{j}}{z - \beta_{j}}$($|\beta_{j}| < 1$),则 $\phi(f) = \sum_{j=1}^{n} c_{j}f(\beta_{j})$.

(2) 若 $k(z) = n!(z-\beta)^{-n-1}$($|\beta| < 1$),则 $\phi(f) = f^{(n)}(\beta)$.

(3) 若 $k(z) = \sum_{j=0}^{n} c_{j}z^{-j-1}$,则 $\phi(f) = \sum_{j=0}^{n} c_{j}a_{j}$.
\end{example}

下面的例子描绘了 (3) 中 $c_{0} = \cdots = c_{n-1} = 0$ 且 $c_{n} = 1$ 时的 $S_{\phi}^{p}$.它在 $\phi$ 为连续函数时描绘 $S_{\phi}^{p}$ 是重要的.

\begin{example}\label{exa:2.16.3}
当 $k(z) = z^{-(n+1)}$ 时,若 $p \neq 1$ 则 $S_{k}^{p} = \{z^{n}\}$,p = 1 时
\begin{equation*}
S_{k}^{1} = \left\{\gamma \prod_{j=1}^{n} (z - \alpha_{j})(1 - \bar{\alpha}_{j}z) \in S^{1}; \gamma > 0 \ \text{且} \ \alpha_{j} \in \mathbb{C}, |\alpha_{j}| \leq 1 \right\}.
\end{equation*}
\end{example}

\begin{proof}
p $\neq$ 1 时由定理~\ref{thm:2.14.4} 显然.p = 1 时,$S_{k}^{1}$ 包含右边由命题~\ref{prop:2.16.1} 显然.证明 $S_{k}^{1}$ 包含于右边.对 $f \in S_{k}^{1}$,令 $F(z) = f(z)/z^{n}$,由命题~\ref{prop:2.16.1},$F(z) \geq 0$($z \in \partial D$),故 $F(z) = \sum_{j=-n}^{n} c_{j}z^{j}$,$\bar{c}_{j} = c_{-j}$ 且 $c_{n} \neq 0$.$F(z)$ 是有理函数,且 $F(z) \geq 0$($z \in \partial D$),由镜像原理,若 $F(\alpha) = 0$ 则 $F(1/\bar{\alpha}) = 0$.$f(z) = z^{n}F(z)$ 是 2n 次多项式,有 2n 个零点.$\partial D$ 上的零点为偶重.于是
\begin{equation*}
f(z) = \gamma \prod_{j=1}^{n} (z - \alpha_{j})(1 - \bar{\alpha}_{j}z),
\end{equation*}
其中 $|\alpha_{j}| \leq 1$ 且 $\gamma > 0$.
\end{proof}

\begin{lemma}\label{lem:2.16.4}
设 $q_{j}$($1 \leq j \leq \ell$)为非常值的内函数,$q = \prod_{j=1}^{\ell} q_{j}$ 且 h 为外函数.设 $K = \overline{zq}|h|/h$,则对任意 $\{b_{j}\}_{j=1}^{\ell} \subset D$,存在 $f \in S_{K}^{1}$,使得 $f(b_{1}) = \cdots = f(b_{\ell}) = 0$.
\end{lemma}

\begin{proof}
对 a.e. $\theta$,
\begin{equation*}
\overline{q(e^{i\theta})}\prod_{j=1}^{\ell} (q_{j}(e^{i\theta}) - q_{j}(0))(1 - \overline{q_{j}(0)}q_{j}(e^{i\theta})) \geq 0,
\end{equation*}
故选取常数 $\gamma > 0$,由命题~\ref{prop:2.16.1},
\begin{equation*}
g = \gamma h\prod_{j=1}^{\ell} (q_{j} - q_{j}(0))(1 - \overline{q_{j}(0)}q_{j})
\end{equation*}
属于 $S_{K}^{1}$.因 $g/z^{\ell} \in H^{1}$,对任意 $\{b_{j}\}_{j=1}^{\ell} \subset D$,存在常数 $\delta > 0$,使得
\begin{equation*}
f = \delta \frac{g}{z^{\ell}}\prod_{j=1}^{\ell} (z - b_{j})(1 - \bar{b}_{j}z)
\end{equation*}
属于 $S_{K}^{1}$,故 f 为所求函数.
\end{proof}

\begin{lemma}\label{lem:2.16.5}
设 $K = \bar{z}^{m+1}|g|/g$,m 为非负整数且 g 为 $S^{1}$ 的露点.任意 $f \in S_{K}^{1}$ 的因式分解中的内函数,是在 D 中零点个数不超过 $(m+1)$ 的有限 Blaschke 积.
\end{lemma}

\begin{proof}
若对 $\{a_{j}\}_{j=1}^{m+1} \subset D$ 存在 $f \in S_{K}^{1}$,使得 $f(a_{1}) = \cdots = f(a_{m+1}) = 0$,则存在 $\gamma > 0$ 与 $h \in H^{1}$,使得
\begin{equation*}
f(z) = \gamma \prod_{j=1}^{m+1} (z - a_{j})(1 - \bar{a}_{j}z) h.
\end{equation*}
于是由例~\ref{exa:2.16.3},存在 $\gamma_{0} > 0$,可令 $\gamma_{0}z^{m+1}h \in S_{K}^{1}$.此时由引理~\ref{lem:2.14.3},a.e. $\theta$ 有 $e^{i\theta}h(e^{i\theta})/g(e^{i\theta}) \geq 0$,由露点的定义,g 是 zh 的常数倍.于是 g 不是外函数,与 g 是 $S^{1}$ 的露点矛盾.若 f = qh 可分解且 q 为奇异内函数,则 $q = \prod_{j=1}^{m+1} q_{j}$ 可分解,其中 $q_{j}$($1 \leq j \leq m+1$)不是常数.于是取 $\delta > 0$,$|\alpha_{j}| \leq 1$($1 \leq j \leq m+1$)且
\begin{equation*}
F(z) = \delta \prod_{j=1}^{m+1} (q_{j} - \alpha_{j})(1 - \bar{\alpha}_{j}q_{j}) h,
\end{equation*}
则 F 可属于 $S_{K}$.选取 $q_{j} - \alpha_{j}$ 使在 D 中有零点,则 F 有 $(m+1)$ 个零点.这与上面所证矛盾,故 f 的内函数是所求的 Blaschke 积.
\end{proof}

\begin{lemma}\label{lem:2.16.6}
设 $K = \overline{z}^{m+1}|g|/g$,m 为非负整数且 g 为 $S^{1}$ 的露点.此时
\begin{equation*}
S_{K}^{1} = (\{\gamma\} \times S_{\bar{z}^{m+1}}^{1} \times g) \cap S^{1}.
\end{equation*}
\end{lemma}

\begin{proof}
对 m 用归纳法证明.令 $K_{m} = \overline{z}^{m+1}|g|/g$.若 m = 0,因 g 是 $S^{1}$ 的露点,$S_{K_{0}}^{1} = \{g\}$.设对 $1 \leq m \leq n$ 本引理成立,即 $S_{K_{m}}^{1} = (\{\gamma\} \times S_{\bar{z}^{m+1}}^{1} \times g) \cap S^{1}$ 成立.$S_{K_{n+1}}^{1} \supseteq (\{\gamma\} \times S_{\bar{z}^{n+2}}^{1} \times g) \cap S^{1}$ 由例~\ref{exa:2.16.3} 显然,故证明 $S_{K_{n+1}}^{1} \subseteq (\{\gamma\} \times S_{\bar{z}^{n+2}}^{1} \times g) \cap S^{1}$.若 $f \in S_{K_{n+1}}^{1}$,则
\begin{equation*}
F = (f + z^{n+1}g)/2 \in S_{K_{n+1}}^{1}.
\end{equation*}
因 F 不是 $S^{1}$ 的端点,由定理~\ref{thm:2.15.3},F 不是外函数.由引理~\ref{lem:2.16.5},存在某个 $1 \leq \ell \leq n+1$,$\{a_{j}\}_{j=1}^{\ell} \subset D$ 与外函数 $h \in H^{1}$,使得
\begin{equation*}
F = \prod_{j=1}^{\ell} \frac{z - a_{j}}{1 - \bar{a}_{j}z} h.
\end{equation*}
对每个 j,因 $\mathrm{Re}(1/(1-\bar{a}_{j}z)) \geq 0$,由定理~\ref{thm:2.15.3} 与例~\ref{exa:2.15.7},$(1-\bar{a}_{j}z)^{-1}$ 是外函数.有限个外函数的乘积仍为外函数,故令
\begin{equation*}
k = \frac{1}{\alpha}\prod_{j=1}^{\ell} \frac{h}{(1 - \bar{a}_{j}z)^{2}},
\end{equation*}
并选取 $\alpha > 0$ 使 $\|k\|_{1} = 1$,则 k 是 $H^{1}$ 的外函数,且
\begin{equation*}
F = \alpha \prod_{j=1}^{\ell} (z - a_{j})(1 - \bar{a}_{j}z) k.
\end{equation*}
\begin{sloppypar}
于是 a.e. $\theta$ 有 $e^{i\ell\theta}k(e^{i\theta})/F(e^{i\theta}) \geq 0$.又因 a.e. $\theta$ 有 $F(e^{i\theta})e^{-i(n+1)\theta}/g(e^{i\theta}) \geq 0$,故 a.e. $\theta$ 有 $e^{i\ell\theta}k(e^{i\theta})e^{-i(n+1)\theta}/g(e^{i\theta}) \geq 0$,即 $k \in S_{K_{n-\ell+1}}^{1}$.
\end{sloppypar}由归纳假设,
\begin{equation*}
k = \beta \prod_{j=1}^{n-\ell+1} (z - b_{j})(1 - \bar{b}_{j}z) g,
\end{equation*}
其中 $\beta > 0$ 且 $\{b_{j}\}_{j=1}^{n-\ell+1} \subset \partial D$.令 $a_{\ell+j} = b_{j}$($1 \leq j \leq n - \ell + 1$),则
\begin{equation*}
F = \alpha\beta \prod_{j=1}^{n+1} (z - a_{j})(1 - \bar{a}_{j}z) g.
\end{equation*}
令
\begin{equation*}
p = 2\alpha\beta \prod_{j=1}^{n+1} (z - a_{j})(1 - \bar{a}_{j}z) - z^{n+1},
\end{equation*}
则 $f = 2F - z^{n+1}g = pg$.因 $f \in S_{K_{n+1}}^{1}$,a.e. $\theta$ 有 $e^{-i(n+1)\theta}g(e^{i\theta})^{-1}f(e^{i\theta}) \geq 0$,故 a.e. $\theta$ 有 $e^{-i(n+1)\theta}p(e^{i\theta}) \geq 0$.于是 $p \in \{\gamma\} \times S_{\bar{z}^{n+2}}^{1}$,故 f 属于 $(\{\gamma\} \times S_{\bar{z}^{n+2}}^{1} \times g) \cap S^{1}$.
\end{proof}

\begin{lemma}\label{lem:2.16.7}
设 $k \notin H^{\infty}$ 且 $S_{k}^{1}$ 是非空 $S^{1}$ 的 *弱紧集.则存在非负整数 m,使得任意 $f \in S_{k}^{1}$ 的因式分解中的内函数,是在 D 中零点个数不超过 m 的有限 Blaschke 积.
\end{lemma}

\begin{proof}
若对 D 中任意有限个点,以之为零点的函数存在于 $S_{k}^{1}$,则由引理~\ref{lem:2.16.4},对每个自然数 n,
\begin{equation*}
f_{n}\left(\frac{1}{2}\right) = f_{n}\left(\frac{1}{3}\right) = \dots = f_{n}\left(\frac{1}{n}\right) = 0
\end{equation*}
的 $f_{n}$ 存在于 $S_{k}^{1}$.因 $S_{k}^{1}$ 是 *弱紧的,$\{f_{n}\}$ 的子列在 $H^{1} = (C/zA)^{*}$(参见定理~\ref{thm:2.13.3})的 *弱拓扑下收敛到 $f \in S_{k}^{1}$.对每个 n,点 $z = 1/n$ 处的点赋值由命题 2.9.4 关于 *弱拓扑连续,故 $f(1/n) = 0$($n = 2, 3, \ldots$).由唯一性定理,$f \equiv 0$.这与 $f \in S_{k}^{1}$ 矛盾.因此,由于引理~\ref{lem:2.16.4} 表明任意 $f \in S_{k}^{1}$ 的因式分解中的内函数是有限 Blaschke 积,本引理证明完毕.
\end{proof}

\begin{theorem}\label{thm:2.16.8}
设 $\phi \in C/zA$,即 $k \in C$.此时存在非负整数 m 与 $S^{1}$ 的露点 g,使得
\begin{equation*}
S_{\phi}^{1} = S_{k}^{1} = \left\{\gamma \prod_{j=1}^{m} (z - a_{j})(1 - \bar{a}_{j}z) g \in S^{1}; \gamma > 0, \{a_{j}\}_{j=1}^{m} \subset D \right\}.
\end{equation*}
\end{theorem}

\begin{proof}
由假设与定理~\ref{thm:2.13.3},$S_{\phi}^{1}$ 是 *弱紧的.由引理~\ref{lem:2.16.7},任意 $f \in S_{\phi}^{1}$ 的因式分解中的内函数是有限 Blaschke 积,其零点个数不超过某个常数.设 m 为其零点最大个数,则存在 $f \in S_{\phi}^{1}$,$f = \prod_{j=1}^{m} \frac{z - a_{j}}{1 - \bar{a}_{j}z} h$,其中 $h \in H^{1}$ 是外函数.
\begin{equation*}
g = \gamma h / \prod_{j=1}^{m} (1 - \bar{a}_{j}z)^{2},
\end{equation*}
选取 $\gamma > 0$ 可使 $g \in S^{1}$.
\begin{equation*}
\frac{f(e^{i\theta})}{e^{im\theta}g(e^{i\theta})} = \frac{1}{\gamma}\frac{\prod_{j=1}^{m} (e^{i\theta} - a_{j})(1 - \bar{a}_{j}e^{i\theta})}{e^{im\theta}} \geq 0 \ \text{a.e.} \ \theta,
\end{equation*}
故由命题~\ref{prop:2.16.1},$z^{m}g \in S_{\phi}^{1}$.一旦知道 g 是 $S^{1}$ 的露点,令 $K = \bar{z}^{m+1}|g|/g$ 并用引理~\ref{lem:2.16.6},定理证明完毕.若 g 不是 $S^{1}$ 的露点,由命题~\ref{prop:2.15.5},存在 $G \in S^{1}$($g \neq G$)使得 $z^{m}G \in S_{\phi}^{1}$.此时 $z^{m}(g + G)/2 \in S_{\phi}^{1}$,但由定理~\ref{thm:2.15.3},$(g + G)/2$ 不是外函数,这与 $S_{\phi}^{1}$ 中函数的内函数部分是 Blaschke 积且零点个数不超过 m 相矛盾.
\end{proof}

\begin{theorem}\label{thm:2.16.9}
设 $1 < p < \infty$,$1/p + 1/q = 1$ 且 $\phi \in (H^{p})^{*}$,$\|\phi\| = 1$.设 K 为 $\phi$ 的极值核,G 为满足 a.e. $\theta$ 有 $|K(e^{i\theta})| = |G(e^{i\theta})|$ 的外函数.设 $\phi_{0} \in (H^{1})^{*}$ 由 $K_{0} = K/G$ 定义.若 $\phi_{0} \in C/zA$,则存在 $\{a_{j}\}_{j=1}^{n} \subset \overline{D}$ 与 $\gamma > 0$,使得
$K_{0} = \bar{z}^{n+1}\frac{|g|}{g}$ 且 g 是 $S^{1}$ 的露点,
\begin{equation*}
G = \gamma^{1/q}\prod_{j=1}^{s} (1 - \bar{a}_{j}z)^{2/q}\left[ \prod_{j=s+1}^{n} (z - a_{j})(1 - \bar{a}_{j}z) \right]^{1/q} g^{1/q},
\end{equation*}
且 $S_{\phi}^{p}$ 由形如
\begin{equation*}
\gamma^{1/p}\prod_{j=1}^{s} (z - a_{j})(1 - \bar{a}_{j}z)^{2/p - 1}\left[ \prod_{j=s+1}^{n} (z - a_{j})(1 - \bar{a}_{j}z) \right]^{1/p} g^{1/p}
\end{equation*}
的 $S^{p}$ 的元素组成.其中 $|a_{j}| < 1$($1 \leq j \leq s$),$|a_{j}| = 1$($s + 1 \leq j \leq n$)且 $\gamma > 0$.
\end{theorem}

\begin{proof}
因 $1 < p < \infty$,由定理~\ref{thm:2.14.4},存在 $F \in H^{p}$ 使得 $S_{\phi}^{p} = \{F\}$.此时因 K 是 $\phi$ 的极值核,由引理~\ref{lem:2.14.3},a.e. $\theta$ 有 $|K(e^{i\theta})| = |F(e^{i\theta})|^{p/q}$.由定理~\ref{thm:2.8.3},F = Qh 可作因式分解.令 G = $h^{q/p}$,则 G 是外函数,且 a.e. $\theta$ 有 $|K(e^{i\theta})| = |G(e^{i\theta})|$,且 F = $QG^{p/q}$.因 $K_{0} = K/G$,利用 q/p + 1 = q,由引理~\ref{lem:2.14.3},a.e. $\theta$ 有 $e^{i\theta}Q(e^{i\theta})G^{q}(e^{i\theta})K_{0}(e^{i\theta}) \geq 0$.因 $|QG^{q}| = |F|^{p}$,$QG^{q} \in S^{1}$,故 $QG^{q} \in S_{\phi_{0}}^{1}$.因 $\phi_{0} \in C/zA$,由定理~\ref{thm:2.16.8},$K_{0} = \overline{z}^{n+1}|g|/g$,g 是 $S^{1}$ 的露点,且
\begin{equation*}
QG^{q} = \gamma \prod_{j=1}^{n} (z - a_{j})(1 - \bar{a}_{j}z) g
\end{equation*}
存在 $\gamma > 0$ 与 $\{a_{j}\}_{j=1}^{n} \subset \overline{D}$ 和 g 满足上式.这里 $\{a_{j}\}_{j=1}^{n}$ 可看作满足 $|a_{j}| < 1$($1 \leq j \leq s$)且 $|a_{j}| = 1$($s + 1 \leq j \leq n$).因此
\begin{equation*}
Q = \prod_{j=1}^{s} \frac{z - a_{j}}{1 - \bar{a}_{j}z},
\end{equation*}
且
\begin{equation*}
G^{q} = \gamma \prod_{j=1}^{s} (1 - \bar{a}_{j}z)^{2}\prod_{j=s+1}^{n} (z - a_{j})(1 - \bar{a}_{j}z) g.
\end{equation*}
因 $F = QG^{q/p}$,
\begin{equation*}
F = \gamma^{1/p}\prod_{j=1}^{s} (z - a_{j})(1 - \bar{a}_{j}z)^{2/p - 1}\left[ \prod_{j=s+1}^{n} (z - a_{j})(1 - \bar{a}_{j}z) \right]^{1/p} g^{1/p}.
\end{equation*}
\end{proof}

定义 $\phi$ 的核 k 为有理函数的情形在应用上很重要,被许多人所研究.例如,参见 Duren 的专著 \cite{ch2:11} 的第 8 章.若 k 为有理函数,则 $k \in C$,故可应用定理~\ref{thm:2.16.8}.另外,当 $1 < p < \infty$ 时,已知 k 的极值核 K 又是有理函数,故可应用定理~\ref{thm:2.16.9}.定理~\ref{thm:2.16.8} 在 \cite{ch2:8} 中,定理~\ref{thm:2.16.9} 在 \cite{ch2:12} 中被证明.定理~\ref{thm:2.16.8} 中 $\phi \notin C/zA$ 时 $S_{\phi}^{1}$ 的研究见 \cite{ch2:13}.
